\section{\enumo by example}
\label{sec:overview}
\enumo is a DSL that provides
 the operators needed to build a
 theory explorer for \eqsat,
 driven by \eqsat, \textit{\`a la carte}.
Users define their own term enumeration strategies and
customize how the resulting rules are processed
and combined.

To introduce the basics of using \enumo,
 we will walk through a simple example
 learning rules over the domain of rational arithmetic.
 \autoref{sec:overview-basic}
 recreates prior work on theory exploration \citep{ruler}
 in a few lines of \enumo code.
\autoref{sec:guided-enumeration}
 goes beyond the capabilities of existing tools
 by employing \enumo's workload construction operators
 to guide term enumeration.
\autoref{subsec:ffoverview}
 demonstrates \enumo's ruleset manipulation primitives,
 including a new way to learn rules
 \emph{without} an interpreter.

\subsection{\enumo basics: Learning rules for rational arithmetic}
\label{sec:overview-basic}
\label{subsec:basic}
Consider the task of
  learning rules for the domain of
 rational arithmetic
 (a grammar with operators \texttt{+, -, $\times$, /, abs, $\sim$})
 for which we assume the user can provide a
 concrete evaluator (a standard recursive
 interpreter).
As with any theory exploration tool,
 the first step is to
 enumerate terms from the domain.
This is typically accomplished by giving a grammar
 to the theory explorer, which then exhaustively enumerates
 terms up to some depth.
In \enumo, however, it is up to the user
 to construct a \emph{workload} that enumerates terms
 by composing various workload primitives and combinators.
The simplest workload is a set of \sexps:
\begin{lstlisting}[name=overview]
  leaves = { a b c -1 0 1 }
  grammar = {
    EXPR                $\label{ln:bare-expr}$
    (- EXPR)
    (abs EXPR)
    (+ EXPR EXPR)
    (- EXPR EXPR)
    (* EXPR EXPR)
    (/ EXPR EXPR)
  }
\end{lstlisting}

Above, the \C{leaves} workload is a set of six atomic \sexps
 that represent the atoms of the domain:
 symbols \C{a}, \C{b}, \C{c}, and constants \C{-1}, \C{0}, and \C{1}.
The \C{grammar} workload is a set of
  \sexps that represent the grammar of the domain.
The \C{EXPR} symbol is intended as a placeholder
 but has no special meaning to \enumo.
By convention, we fully capitalize these placeholder
  symbols to differentiate them from symbols
  and operators in the domain.

The \C{plug} operator allows the user to compose two workloads
 $\W_1$ and $\W_2$ by replacing occurrences of a given symbol $x$
 in $\W_1$ with \sexps from $\W_2$.
Using \C{plug},
 the user can construct a workload that exhaustively enumerates terms
 up to a particular depth:

\begin{lstlisting}[name=overview]
  rationals_depth1 = grammar.plug("EXPR", leaves)           $\label{ln:plug1}$
  rationals_depth2 = grammar.plug("EXPR", rationals_depth1) $\label{ln:plug2}$
\end{lstlisting}

Note that \C{$\W_1$.plug("x", $\W_2$)} yields
 all possible combinations of replacing symbol \C{"x"} in $\W_1$ with an \sexp from $\W_2$
 (see \autoref{sec:slide} for detailed semantics).
Because the \C{grammar} on line \ref{ln:bare-expr} includes \C{EXPR},
\C{grammar.plug("EXPR", W)} will include all \sexps in \C{W}.
Lines \ref{ln:plug1} and \ref{ln:plug2}
 will enumerate all terms up to depth 1 and 2, respectively.

With a workload in hand
 that represents enumerated terms from the domain,
 we can proceed with writing an \enumo program to learn rules.
Initially,
 we will learn rules following
 the conventional approach
 of \citet{ruler} and \citet{sat19};
 in later subsections, we will show a more advanced \enumo program.
First, we learn rules of depth 1 (D1):

\begin{lstlisting}[name=overview]
  candidates1 = rationals_depth1
                 .to_egraph()                       $\label{ln:to-egraph}$
                 .find_candidates()                     $\label{ln:cvec-match}$
  (valid1, _) = candidates1.partition(|c| c.is_valid()) $\label{ln:validate}$
  rules1 = valid1.minimize([]) $\label{ln:minimize}$
\end{lstlisting}

Line $\ref{ln:to-egraph}$ converts the D1 workload to an \egraph
 by evaluating the workload (according to the semantics in \autoref{sec:slide})
 and adding the resulting \sexps to the \egraph.
Unlike the untyped \sexps in the workload,
 the \egraph is typed according to the domain $L$
 (in this case, rational arithmetic).\footnote{
  An exception occurs if \sexps cannot be parsed into the syntax for domain $L$.
 }
Once the \egraph is constructed,
 we use \C{find\_candidates}
 to generate rule candidates (line $\ref{ln:cvec-match}$).
This method uses the \emph{characteristic vector} (\cvec) matching
 approach from \citet{ruler},
 evaluating the terms in the \egraph
 with the user-provided interpreter on a sampling of constants.
Terms that evaluate identically on all inputs
 are likely to be \textit{equivalent}, and are thus candidates for rules.
Candidates are not necessarily sound,
 so we must validate them on line $\ref{ln:validate}$
 using a user-provided verifier
 (over the rationals, we use Z3~\citep{z3}).

Finally, on line $\ref{ln:minimize}$,
 we minimize the candidates
 to eliminate redundant rules.
The \C{minimize} operator in \enumo
 is parameterized over a scheduling \C{Strategy}.
With a \C{Strategy}, a user can control
 how much redundancy is permissible in the ruleset,
 measured via ``derivability'', as defined
 in \autoref{subsec:derivability}.
Minimize works as follows:
rules from \C{valid1} are added one by one
to a new, initially empty ruleset \C{rules1}. A given rule $r \in$ \C{valid1}
is added to \C{rules1} if the rules currently in \C{rules1} \textit{cannot}
derive $r$ within a given number of \eqsat iterations.

At this point,
 this small \enumo program has emulated the
 behavior of \citet{ruler}'s \ruler tool
 to learn D1 rules.
We can now learn rules of depth 2 (D2)
 by repeating the process:
\pagebreak

\begin{lstlisting}[name=overview]
  candidates2 = rationals_depth2 # from line (*@\ref{ln:plug2}@*)
                 .to_egraph()
                 .compress(rules1) $\label{ln:compress}$
                 .find_candidates()
  rules2 = candidates2.minimize(rules1)
  all_rules = rules2.union(rules1)
\end{lstlisting}

Again following \citet{ruler},
 learning the D2 rules benefits from the D1 rules
 in two ways.
First,
 on line $\ref{ln:compress}$,
 we
 compress the D2 \egraph
 prior to finding candidates.
This not only shrinks the \egraph
 and makes finding candidates more efficient,
 but also prevents learning candidates that are implied by the D1 rules.
Similarly,
 we pass the D1 rules to the \C{minimize} operator
 to minimize candidates not only with respect to each other,
 but also with respect to the D1 rules.
 Finally,
 we compose \C{rules2} and \C{rules1} into \C{all_rules}.

\subsection{Guided enumeration to find ``deeper'' rules}
\label{sec:guided-enumeration}

Notice that we have just implemented the entire \ruler tool in about twenty
  lines of \enumo.
\enumo workload operators can easily express exhaustive term enumeration, and
  the candidate generation and selection techniques used in \ruler are also
  supported in \enumo.
However, while \ruler can only perform a few iterations before getting stuck
  due to the exponential growth of enumerated terms, \enumo programs can
  express subsets of the entire term space, making it easy to scale beyond what
  is possible with exhaustive enumeration.

In the previous example (\autoref{sec:overview-basic}), we learned rules over
the \C{rational} domain, which is complicated by rewrites involving division (/)
that are only conditionally true---for example, the rule \C{(/ a a) \rewritesto\,1}
holds only when \C{a} is nonzero. To address this problem, we implemented
a version of the \C{rational} domain in \enumo that supports conditional statements.

Suppose we want to find the rule
  \C{(/ a a) \rewritesto\, (if a 1 (/ a a))},
  which is a version of the rule
  \C{(/ a a) \rewritesto\, 1} that inserts
  a condition so that the term simplifies to 1
  only when the denominator is nonzero.
The right-hand side of this rule is a large term, so
  exhaustively enumerating all terms from the domain up to its size will
  result in an \egraph with over 3 million terms, far too many
  to be feasible for rule inference.
In \ruler, there is no simple mechanism by which a domain expert can
  narrow the search space in order to learn deeper rules; however,
  \enumo enables users to leverage their domain expertise to find better, deeper
  rules than is possible with exhaustive term enumeration.

Our key goal in this section is to learn conditional versions of unsound rules,
  so let's first find the unsound rule candidates from \autoref{sec:overview-basic}:

\begin{lstlisting}[name=guided]
  # rule candidates over terms up to depth 2
  all_candidates = candidates1.union(candidates2)
  # partition rule candidates using domain-provided rule validator
  (sound, unsound) = all_candidates.partition(|rule| rule.is_valid())
\end{lstlisting}

Now, we can construct a workload that adds a check for
  division by zero:

\pagebreak

\begin{lstlisting}[name=guided]
  guard_wkld = Workload.new()
  guard_pattern = { (if GUARD THEN ELSE) } $\label{ln:guard-pat}$
  for rule in unsound:
    # domain-specific function that returns a workload consisting of terms that
    # appear as the second argument to division
    denominators = rule.denominators() $\label{ln:ext}$
    # construct terms that match the unsound rule candidate, but with a check
    # for division by zero
    guard_wkld = guard_wkld.union(
      guard_pattern
        .plug("GUARD", denominators)
        .plug("THEN", rule.rhs)
        .plug("ELSE", rule.lhs)
    )

\end{lstlisting}

Above, \C{guard_pattern} (line $\ref{ln:guard-pat}$) is a workload consisting of
  a single \sexp, which will serve as a pattern for the workload we are
  constructing.
Note that because \enumo is an embedded DSL, it is possible to encode
  domain-specific extensions simply by writing custom functions.
For example, on line $\ref{ln:ext}$ above, we use a domain-specific function
  to construct a workload.
We loop over the unsound rules, incrementally building the workload
  by using the \C{plug} operator to make a guarded version of each unsound rule.
With this workload in hand, we are ready to learn rules:

\begin{lstlisting}[name=guided]
  candidates = guard_wkld.to_egraph().compress(rules2).find_candidates()
  guard_rules = candidates.minimize(rules2)
\end{lstlisting}

This step closely mirrors the process of learning D2 rules in the
  previous subsection: we convert the workload to an \egraph, compress
  the \egraph using the rules we've already learned, and find candidates.
Finally, we minimize the candidates again using the existing rules.
The final ruleset contains the rules
  \C{(/ a a) \rewritesto \xspace (if a 1 (/ a a))} and
  \C{(/ 0 a) \rewritesto \xspace (if a 0 (/ 0 a))}, both of which are useful, sound
  rewrite rules that avoid unsoundness when the denominator could be zero.
Importantly, we find these rules without enumerating all depth-3
  terms.

\subsection{Learning refined rules with ruleset manipulation}
\label{subsec:ffoverview}

Now, let's consider an alternate approach to candidate generation.
Suppose we want to learn rules for
  transcendental functions (e.g., trigonometric operators).
For the examples in \autoref{sec:overview-basic}
and \autoref{sec:guided-enumeration}, we used cvec matching,
  which requires the user to
  implement an interpreter for the domain.
For transcendental functions,
  equality is undecidable~\citep{boehm}, so
  cvec matching is not possible.
However, these functions \textit{can} be represented
  in terms of other functions
  over rational and complex
  domains\footnote{\url{https://en.wikipedia.org/wiki/Cis_(mathematics)}},
  for which there exist various identities.
For example,
  the functions sine and cosine can
  be represented mathematically by
  $\sin(x) = \frac{\cis(x) - \cis(-x)}{2i}$ and
  $\cos(x) = \frac{\cis(x) + \cis(-x)}{2}$
  where $\cis(x)$ is $e^{ix}$.

Leveraging the compositional nature of \enumo operators,
  we can first synthesize rewrite rules
  over rationals and then
  use them to learn rules for
  trigonometric functions without needing
  to evaluate trigonometric terms directly.
We begin with a set of rewrite rules over rationals
  synthesized previously using an \enumo program.
\pagebreak
\begin{lstlisting}[name=trig]
  initial_rules = Ruleset.from_file("initial.rules")
\end{lstlisting}
Next, we define \emph{exploratory} rules (\autoref{sec:lift})
  that express the
  trigonometric operators in terms of
  the rational and complex operators.
These rules are typically handwritten:
\begin{lstlisting}[name=trig]
  explore = [
    "(sin a) ==> (/ (- (cis a) (cis (- a))) (* 2 I))",
    "(cos a) ==> (/ (+ (cis a) (cis (- a))) 2)",
    "(tan a) ==> (* I (/ (- (cis (- a)) (cis a)) (+ (cis (- a)) (cis a))))"
  ]
\end{lstlisting}
Now, we construct a workload representing trigonometric terms:
\begin{lstlisting}[name=trig]
  consts = { 0 (/ PI 6) (/ PI 4) (/ PI 3) (/ PI 2) PI (* PI 2) }
  wkld = { (OP VAL) }.plug("OP", { sin cos tan }).plug("VAL", consts)
          .filter(Filter.Not(Filter.Contains("(tan (/ PI 2))")))
\end{lstlisting}
This workload represents terms that have
  a single trigonometric operator applied
  to a constant value; notice that we can easily filter out
  \C{(tan (/ PI 2))}, which is undefined.
Finally, we run the \ff algorithm on this workload to discover
  new rules from our initial and exploratory rules.

The \ff algorithm is explained in detail in \autoref{sec:lift}, but at a high level,
it identifies candidates by running known rewrite
  rules and considering merged \eclasses as rule candidates.
By definition, all rules found using this algorithm are derivable from the
  starting ruleset, but the rulesets generated can still be valuable in practice.
In this case, \ff allows us to find rules over the trigonometric operators
  directly, rather than needing to rewrite through large complex terms.

